% Gabarit pandoc des corrigés publiés (brut et guidé). Contenu sans marque ni nom de site.
% Compilation : lualatex (latexmk -lualatex).
% Le Markdown reste la source unique : ne pas éditer ce fichier .tex à la main.
\documentclass[11pt,a4paper]{article}

% --- Moteur, polices (Latin Modern) et langue ---------------------------------
\usepackage{iftex}
\ifLuaTeX\else\errmessage{Ce document se compile avec lualatex}\fi
\usepackage{amsmath,amssymb,mathtools}
\usepackage{fontspec}
% mathrm=sym : \mathrm prend ses glyphes dans la police mathématique (ℓ, μ, Ω), comme KaTeX.
\usepackage[math-style=TeX,bold-style=TeX,mathrm=sym]{unicode-math}
% Repli pour les glyphes absents de Latin Modern (indices et exposants Unicode,
% flèches, filets des schémas en texte, symboles) : DejaVu, puis Latin Modern Math
% (ℓ, lettres mathématiques Unicode des légendes).
\directlua{
  luaotfload.add_fallback("repliTexte", {"DejaVuSerif:mode=harf;", "DejaVuSans:mode=harf;", "file:latinmodern-math.otf:mode=harf;"})
  luaotfload.add_fallback("repliMono", {"DejaVuSansMono:mode=harf;", "DejaVuSans:mode=harf;", "file:latinmodern-math.otf:mode=harf;"})
}
\setmainfont{Latin Modern Roman}[RawFeature={fallback=repliTexte}]
\setsansfont{Latin Modern Sans}[RawFeature={fallback=repliTexte}]
\setmonofont{Latin Modern Mono}[RawFeature={fallback=repliMono}]
\setmathfont{Latin Modern Math}
\usepackage[french]{babel}
\frenchsetup{SmallCapsFigTabCaptions=false}

% --- Mise en page --------------------------------------------------------------
\usepackage[a4paper,margin=2.2cm,headheight=14pt,footskip=1.2cm]{geometry}
\usepackage{microtype}
\setlength{\emergencystretch}{4em}
\tolerance=1500
\setlength{\parindent}{0pt}
\setlength{\parskip}{0.55\baselineskip plus 2pt minus 1pt}
\usepackage{xcolor}
% Gris nommés : babel-french rend « ! » actif dans le document.
\definecolor{gris2}{gray}{0.98}
\definecolor{gris3}{gray}{0.97}
\definecolor{gris5}{gray}{0.95}
\definecolor{gris35}{gray}{0.65}
\definecolor{gris40}{gray}{0.60}
\definecolor{gris45}{gray}{0.55}
\definecolor{gris50}{gray}{0.50}
\definecolor{gris55}{gray}{0.45}
\definecolor{gris60}{gray}{0.40}
\definecolor{gris65}{gray}{0.35}
\definecolor{gris70}{gray}{0.30}
\definecolor{gris80}{gray}{0.20}
\definecolor{gris85}{gray}{0.15}
\usepackage{graphicx}
\usepackage{float}
\usepackage{caption}
\captionsetup{font=small,labelformat=empty,justification=centering}
\usepackage{booktabs,longtable,array,multirow,calc}
\usepackage{siunitx}
\sisetup{locale=FR}
\usepackage[version=4]{mhchem}
\usepackage{listings}
\usepackage{fancyhdr}
\usepackage{lastpage}
\usepackage[most]{tcolorbox}
\usepackage{enumitem}
\setlist{nosep,leftmargin=*,topsep=2pt}
\usepackage[hidelinks,unicode]{hyperref}
\usepackage{bookmark}
\hypersetup{pdftitle={Centrale-Supélec MP-MPI-PC-PSI — Anglais 2026},pdfauthor={Rédigé par IA},pdfcreator={LaTeX},pdfsubject={Corrigé guidé}}
\usepackage{xurl}

% --- Pandoc : listes, tableaux, figures ----------------------------------------
\providecommand{\passthrough}[1]{#1}
\providecommand{\tightlist}{\setlength{\itemsep}{0pt}\setlength{\parskip}{0pt}}

\makeatletter
\def\maxwidth{\ifdim\Gin@nat@width>\linewidth\linewidth\else\Gin@nat@width\fi}
\def\maxheight{\ifdim\Gin@nat@height>0.55\textheight 0.55\textheight\else\Gin@nat@height\fi}
\makeatother
\setkeys{Gin}{width=\maxwidth,height=\maxheight,keepaspectratio}
\setcounter{secnumdepth}{-\maxdimen}
\providecommand{\pandocbounded}[1]{#1}
\usepackage{longtable,booktabs,array}

% --- Code (listings) -----------------------------------------------------------
\lstset{
  basicstyle=\ttfamily\small,
  columns=fullflexible,
  keepspaces=true,
  breaklines=true,
  breakatwhitespace=false,
  postbreak=\mbox{\textcolor{gris50}{$\hookrightarrow$}\space},
  frame=single,
  rulecolor=\color{gris35},
  framesep=4pt,
  xleftmargin=5pt,xrightmargin=5pt,
  backgroundcolor=\color{gris3},
  keywordstyle=\bfseries,
  commentstyle=\color{gris65},
  stringstyle=\color{gris80},
  showstringspaces=false,
  tabsize=4,
  extendedchars=true,
  aboveskip=6pt,belowskip=6pt,
}

% --- Macros des corrigés (KaTeX ↔ LaTeX) ---------------------------------------
% Les corrigés sont écrits pour KaTeX : on définit ici ce que KaTeX connaît et
% que LaTeX n'a pas, sans jamais retoucher le Markdown.
\providecommand{\R}{\mathbb{R}}
\providecommand{\N}{\mathbb{N}}
\providecommand{\Z}{\mathbb{Z}}
\providecommand{\Q}{\mathbb{Q}}
\providecommand{\C}{\mathbb{C}}
\providecommand{\dd}{\mathrm{d}}
\providecommand{\llbracket}{\lBrack}
\providecommand{\rrbracket}{\rBrack}
\providecommand{\lBrace}{\lbrace}
\providecommand{\rBrace}{\rbrace}
\providecommand{\Box}{\square}
\providecommand{\Diamond}{\lozenge}
\providecommand{\varnothing}{\emptyset}
\providecommand{\argmax}{\operatorname*{arg\,max}}
\providecommand{\argmin}{\operatorname*{arg\,min}}
\renewcommand{\leq}{\leqslant}
\renewcommand{\geq}{\geqslant}

% Équation hors texte : réduite à la largeur de la ligne si elle déborde.
\newsavebox{\boiteequation}
\newcommand{\equationajustee}[1]{%
  \sbox{\boiteequation}{$\displaystyle #1$}%
  \ifdim\wd\boiteequation>\linewidth
    \[\resizebox{\linewidth}{!}{\usebox{\boiteequation}}\]%
  \else
    \[#1\]%
  \fi}

% --- Encadrés pédagogiques (sobres, lisibles en noir et blanc) -----------------
\tcbset{
  encadre/.style={
    enhanced,breakable,sharp corners,
    colback=white,colframe=gris55,
    boxrule=0.4pt,leftrule=2.2pt,
    left=6pt,right=6pt,top=3pt,bottom=3pt,
    before skip=6pt,after skip=6pt,
    fonttitle=\bfseries\small,coltitle=black,
    attach title to upper,after title={\quad},
    parbox=false,
  }
}
\newtcolorbox{indication}[1]{encadre,colframe=gris45,colback=gris2,title={Indication~#1}}
\newtcolorbox{methode}{encadre,colframe=gris70,title={Méthode}}
\newtcolorbox{piege}{encadre,colframe=gris85,boxrule=0.8pt,leftrule=2.2pt,title={Piège}}
\newtcolorbox{jury}{encadre,colframe=gris60,colback=gris5,title={Jury}}
\newtcolorbox{encadrenote}{encadre,colframe=gris40}

% --- Titres de section ---------------------------------------------------------
\usepackage{titlesec}
\titleformat{\section}{\Large\bfseries}{}{0pt}{}
\titleformat{\subsection}{\large\bfseries}{}{0pt}{}
\titleformat{\subsubsection}{\normalsize\bfseries}{}{0pt}{}
\titlespacing*{\section}{0pt}{14pt plus 4pt minus 2pt}{6pt}
\titlespacing*{\subsection}{0pt}{10pt plus 3pt minus 2pt}{4pt}
\titlespacing*{\subsubsection}{0pt}{8pt plus 2pt minus 2pt}{2pt}

% --- En-tête et pied de page ---------------------------------------------------
\pagestyle{fancy}
\fancyhf{}
\fancyhead[L]{\small\textsc{Corrigé guidé} · Centrale-Supélec MP-MPI-PC-PSI}
\fancyhead[R]{\small Anglais 2026}
\fancyfoot[L]{\footnotesize Corrigé rédigé par IA, non officiel}
\fancyfoot[R]{\footnotesize \thepage\,/\,\pageref*{LastPage}}
\renewcommand{\headrulewidth}{0.4pt}
\renewcommand{\footrulewidth}{0pt}
\fancypagestyle{plain}{\fancyhf{}\fancyfoot[R]{\footnotesize \thepage\,/\,\pageref*{LastPage}}\fancyfoot[L]{\footnotesize Corrigé rédigé par IA, non officiel}\renewcommand{\headrulewidth}{0pt}}

\begin{document}
\thispagestyle{plain}
\begin{center}
  {\small\scshape Corrigé guidé}\par\smallskip
  {\LARGE\bfseries Centrale-Supélec MP-MPI-PC-PSI — Anglais 2026\par}\medskip
  {\large Centrale-Supélec · MP-MPI-PC-PSI · Anglais · 2026\par}
\end{center}
\vspace{2pt}\hrule\vspace{8pt}

\emph{Corrigé rédigé par IA, non officiel. Signalez toute erreur.}

\hypertarget{pruxe9sentation-du-sujet}{%
\section{Présentation du sujet}\label{pruxe9sentation-du-sujet}}

\textbf{Format.} L'épreuve dure 4 heures ; calculatrice et tout système électronique sont interdits, et aucun dictionnaire n'est mentionné comme autorisé. Elle comporte un seul exercice :

\begin{longtable}[]{@{}
  >{\raggedright\arraybackslash}p{(\columnwidth - 6\tabcolsep) * \real{0.1101}}
  >{\raggedright\arraybackslash}p{(\columnwidth - 6\tabcolsep) * \real{0.2431}}
  >{\raggedright\arraybackslash}p{(\columnwidth - 6\tabcolsep) * \real{0.2752}}
  >{\raggedright\arraybackslash}p{(\columnwidth - 6\tabcolsep) * \real{0.3416}}@{}}
\toprule\noalign{}
\begin{minipage}[b]{\linewidth}\raggedright
Exercice
\end{minipage} & \begin{minipage}[b]{\linewidth}\raggedright
Nature
\end{minipage} & \begin{minipage}[b]{\linewidth}\raggedright
Longueur
\end{minipage} & \begin{minipage}[b]{\linewidth}\raggedright
Contraintes explicites du sujet
\end{minipage} \\
\midrule\noalign{}
\endhead
\bottomrule\noalign{}
\endlastfoot
Synthèse & synthèse en anglais de quatre documents & 500 mots ± 10 \%, soit 450 à 550 mots, titre inclus & titre obligatoire ; nombre de mots indiqué avec précision à la fin du travail \\
\end{longtable}

Le sujet précise que l'ordre des documents est arbitraire et ne revêt aucune signification particulière. Le barème n'est pas publié dans le sujet. Le fichier du sujet contient aussi, après les quatre pages d'anglais, le sujet « Langue vivante » des autres langues (allemand, arabe, espagnol, italien, portugais, russe, chinois) : ce corrigé ne traite que l'anglais, conformément à l'intitulé de l'épreuve.

\textbf{Pas de rapport de jury.} Centrale-Supélec n'a pas publié de rapport de jury pour la session 2026 au moment de la rédaction. Ce corrigé ne cite donc aucune attente du jury : les conseils qui suivent sont des \textbf{conseils méthodologiques généraux} sur l'exercice de synthèse, présentés comme tels, et non des exigences propres aux correcteurs de cette session.

\textbf{Le dossier.}

\begin{longtable}[]{@{}
  >{\raggedright\arraybackslash}p{(\columnwidth - 6\tabcolsep) * \real{0.1194}}
  >{\raggedright\arraybackslash}p{(\columnwidth - 6\tabcolsep) * \real{0.1286}}
  >{\raggedright\arraybackslash}p{(\columnwidth - 6\tabcolsep) * \real{0.5265}}
  >{\raggedright\arraybackslash}p{(\columnwidth - 6\tabcolsep) * \real{0.1955}}@{}}
\toprule\noalign{}
\begin{minipage}[b]{\linewidth}\raggedright
Document
\end{minipage} & \begin{minipage}[b]{\linewidth}\raggedright
Nature, date
\end{minipage} & \begin{minipage}[b]{\linewidth}\raggedright
Contenu
\end{minipage} & \begin{minipage}[b]{\linewidth}\raggedright
Position
\end{minipage} \\
\midrule\noalign{}
\endhead
\bottomrule\noalign{}
\endlastfoot
D. Dawson & photographie, 2001 & Lucian Freud, vu de dos, palette à la main, peint Elizabeth II, assise, souriante, coiffée d'un diadème ; sur la petite toile, le visage paraît plus lourd et plus sévère. Le portrait est aujourd'hui à la National Portrait Gallery de Londres. & montre le portrait comme une rencontre entre un modèle et un regard \\
D. Smith, \emph{The Guardian} & article, 31/12/2024 & Kim Sajet, directrice de la National Portrait Gallery de la Smithsonian (Washington) : « All portraiture is a lie » ; le portrait, fabrication à quatre (modèle, artiste, commanditaire, public) ; cartels de 140 mots fondés sur les faits ; volonté de rester « above the fray » ; portrait officiel de Trump achevé mais en réserve jusqu'en 2029 ; Picasso : l'art est un mensonge qui fait voir la vérité ; directrice née au Nigeria, élevée en Australie, de nationalité néerlandaise & institution qui revendique la neutralité \\
K. Grovier, \emph{BBC.com} & article, 01/08/2025 & la statue de la Liberté, symbole composite et contesté dès 1886 (suffragettes, conservateurs hostiles aux migrants) ; \emph{Trans Forming Liberty} d'Amy Sherald (la Liberté en femme noire transgenre) ; avertie que l'œuvre pourrait déplaire au président, Sherald annule son exposition à la National Portrait Gallery de la Smithsonian en dénonçant « a culture of censorship » ; décret de mars 2025, rôle de JD Vance ; l'exclusion amplifie l'œuvre : « What's more visible than something hidden? » & critique de la pression politique \\
R. Kimball, \emph{The Spectator} & article d'opinion, 13/08/2025 & éloge du décret « Restoring Truth and Sanity to American History » ; la Smithsonian, fondée en 1846 pour la diffusion du savoir, serait tombée sous une idéologie « divisive, race-centered » ; revue interne des expositions, des textes et des procédures des conservateurs ; 250e anniversaire ; il faut « commandeer » les institutions de la culture d'élite & défense de l'intervention politique \\
\end{longtable}

\textbf{Le thème.} Les quatre documents parlent tous de portraits ou de symboles exposés dans des institutions nationales, et trois d'entre eux de la même institution, la Smithsonian. Ils ne s'accordent pas sur la neutralité des musées : Sajet la revendique, Kimball la nie et veut corriger une dérive idéologique, Grovier y voit une censure. La photographie apporte un contrepoint britannique, plus ancien et apaisé, qui rend visible la thèse de Sajet : un portrait est un regard posé sur un modèle, pas un reflet.

\textbf{Ce qu'il ne faut pas laisser passer} (conseils généraux) :

\begin{itemize}
\tightlist
\item
  les deux contraintes explicites : un \textbf{titre} et le \textbf{nombre de mots} à la fin, dans la fourchette 450-550 ;
\item
  les \textbf{quatre documents}, photographie comprise : un document oublié déséquilibre la synthèse ;
\item
  une synthèse \textbf{confronte} les documents dans chaque partie, au lieu de les résumer l'un après l'autre ;
\item
  la synthèse rend compte du dossier : pas d'opinion personnelle, pas d'informations extérieures, même exactes ;
\item
  les positions des auteurs sont attribuées (\emph{Kimball argues\ldots{}}, \emph{Grovier replies\ldots{}}) : le point de vue partisan de Kimball ne doit pas être repris à son compte.
\end{itemize}

\textbf{Gestion du temps} (proposition) : 1 h 15 de lecture et d'analyse, 30 min pour la confrontation et le plan, 1 h 45 de rédaction au brouillon puis au propre, 30 min de décompte et de relecture.

\hypertarget{partie-i-synthuxe8se-de-documents}{%
\section{Partie I --- Synthèse de documents}\label{partie-i-synthuxe8se-de-documents}}

L'unique exercice est découpé en six étapes, de la lecture au décompte final. La solution complète est donnée à l'étape 5.

\hypertarget{uxe9tape-1-lire-le-dossier-et-analyser-chaque-document}{%
\subsection{Étape 1 --- Lire le dossier et analyser chaque document}\label{uxe9tape-1-lire-le-dossier-et-analyser-chaque-document}}

\begin{indication}{1}

Pour chaque document, notez la nature (photo, reportage, article d'opinion), la date, la source et la thèse en une phrase. Le ton compte : \emph{The Spectator} ne parle pas comme la BBC.

\end{indication}

\begin{indication}{2}

La photographie est un document à part entière. Décrivez ce qu'elle montre (qui regarde qui, où est la toile, ce qu'on voit du visage peint) avant de lui chercher un sens.

\end{indication}

\begin{indication}{3}

Repérez les liens matériels entre les articles : deux d'entre eux parlent du même musée de Washington, et deux citent le même décret présidentiel.

\end{indication}

\hypertarget{solution}{%
\subsubsection{Solution}\label{solution}}

Le tableau de la présentation donne l'analyse de chaque document. Les liens à exploiter sont les suivants :

\begin{longtable}[]{@{}
  >{\raggedright\arraybackslash}p{(\columnwidth - 4\tabcolsep) * \real{0.3165}}
  >{\raggedright\arraybackslash}p{(\columnwidth - 4\tabcolsep) * \real{0.3328}}
  >{\raggedright\arraybackslash}p{(\columnwidth - 4\tabcolsep) * \real{0.3207}}@{}}
\toprule\noalign{}
\begin{minipage}[b]{\linewidth}\raggedright
Lien
\end{minipage} & \begin{minipage}[b]{\linewidth}\raggedright
Documents
\end{minipage} & \begin{minipage}[b]{\linewidth}\raggedright
Usage dans la synthèse
\end{minipage} \\
\midrule\noalign{}
\endhead
\bottomrule\noalign{}
\endlastfoot
même musée : la National Portrait Gallery de la Smithsonian & Smith, Grovier & la neutralité revendiquée par Sajet est mise à l'épreuve par l'affaire Sherald \\
même décret de mars 2025, « Restoring Truth and Sanity to American History » & Grovier, Kimball & lu comme une censure par l'un, comme une restauration de la vérité par l'autre \\
le portrait comme construction & Dawson, Smith, Grovier & partie I \\
le mot \emph{truth} & Smith (Picasso), Kimball (« truth and sanity »), Grovier (« straightforward truth ») & fil directeur de la problématique \\
\end{longtable}

\begin{piege}

Il y a deux National Portrait Galleries : celle de Londres (photographie) et celle de la Smithsonian à Washington (Smith, Grovier). Les confondre fausse toute la synthèse.

\end{piege}

\hypertarget{uxe9tape-2-confronter-les-documents-et-formuler-la-probluxe9matique}{%
\subsection{Étape 2 --- Confronter les documents et formuler la problématique}\label{uxe9tape-2-confronter-les-documents-et-formuler-la-probluxe9matique}}

\begin{indication}{1}

Cherchez une question à laquelle les quatre documents répondent, chacun à sa façon. Le mot qui revient chez Smith, Grovier et Kimball est \emph{truth}.

\end{indication}

\begin{indication}{2}

La problématique doit faire apparaître la tension entre la vérité d'une image et le pouvoir qui l'expose ou la retire.

\end{indication}

\hypertarget{solution-1}{%
\subsubsection{Solution}\label{solution-1}}

\textbf{Problématique} : portraits et symboles nationaux peuvent-ils encore dire la vérité d'une nation lorsque le pouvoir politique prétend en fixer le sens ?

Elle couvre les quatre documents : la photographie et Sajet montrent que le portrait est une construction ; Grovier et Kimball s'opposent sur le rôle de l'État ; Sajet et Grovier doutent qu'on puisse imposer un sens.

\begin{methode}

Une bonne problématique de synthèse se construit à partir des désaccords du dossier, pas de ses seuls points communs. Ici, le désaccord porte sur qui a le droit de dire la vérité d'une image.

\end{methode}

\hypertarget{uxe9tape-3-construire-le-plan}{%
\subsection{Étape 3 --- Construire le plan}\label{uxe9tape-3-construire-le-plan}}

\begin{indication}{1}

Trois parties suffisent pour 500 mots. Chacune doit mobiliser au moins deux documents.

\end{indication}

\begin{indication}{2}

Allez du constat (qu'est-ce qu'un portrait ?) au conflit (qui en contrôle le sens ?), puis à la limite de ce contrôle.

\end{indication}

\hypertarget{solution-2}{%
\subsubsection{Solution}\label{solution-2}}

\begin{longtable}[]{@{}
  >{\raggedright\arraybackslash}p{(\columnwidth - 4\tabcolsep) * \real{0.0986}}
  >{\raggedright\arraybackslash}p{(\columnwidth - 4\tabcolsep) * \real{0.3557}}
  >{\raggedright\arraybackslash}p{(\columnwidth - 4\tabcolsep) * \real{0.5157}}@{}}
\toprule\noalign{}
\begin{minipage}[b]{\linewidth}\raggedright
Partie
\end{minipage} & \begin{minipage}[b]{\linewidth}\raggedright
Idée directrice
\end{minipage} & \begin{minipage}[b]{\linewidth}\raggedright
Documents mobilisés
\end{minipage} \\
\midrule\noalign{}
\endhead
\bottomrule\noalign{}
\endlastfoot
I & Un portrait est une construction, pas un miroir ; les symboles nationaux aussi & Smith (fabrication à quatre), Dawson (le peintre et son modèle), Grovier (statue composite, contestée dès 1886 ; Sherald) \\
II & Les institutions qui exposent ces images deviennent des cibles politiques & Smith (neutralité revendiquée, cartel de Trump), Kimball (décret, revue interne, 250e anniversaire), Grovier (annulation de l'exposition, « culture of censorship ») \\
III & Le sens d'une image se laisse-t-il contrôler ? & Kimball (« commandeer » les institutions), Grovier (l'exclusion amplifie l'œuvre), Smith (le temps long, le portrait en réserve jusqu'en 2029, le regard extérieur) \\
\end{longtable}

Répartition visée : introduction ≈ 75 mots, parties ≈ 130, 145 et 110 mots, conclusion ≈ 50 mots.

\begin{piege}

Le plan « un document par paragraphe » n'est pas une synthèse : c'est une suite de résumés. Chaque partie ci-dessus croise au moins trois sources.

\end{piege}

\hypertarget{uxe9tape-4-choisir-le-titre-et-ruxe9diger-lintroduction}{%
\subsection{Étape 4 --- Choisir le titre et rédiger l'introduction}\label{uxe9tape-4-choisir-le-titre-et-ruxe9diger-lintroduction}}

\begin{indication}{1}

Le titre est obligatoire et compté. Il doit annoncer le sujet et, si possible, la tension (vérité et pouvoir).

\end{indication}

\begin{indication}{2}

L'introduction présente brièvement les quatre documents (auteur ou source, date) et se termine par la problématique, en anglais.

\end{indication}

\hypertarget{solution-3}{%
\subsubsection{Solution}\label{solution-3}}

Titre retenu : \textbf{Painting the Nation: Truth, Power and the Portrait} (8 mots). Autres possibilités : \emph{Who Holds the Brush?} ; \emph{Whose Truth on the Walls?}

Who decides how a nation sees itself? A 2001 photograph by David Dawson shows Lucian Freud painting Queen Elizabeth II, a portrait now held by London's National Portrait Gallery. Three articles, from The Guardian, the BBC and The Spectator, move the question to Washington, where the Smithsonian's museums have become a battlefield during Donald Trump's second term. They ask whether portraits and national symbols can still tell the truth once politics claims them.

(73 mots.)

\begin{methode}

Une phrase d'accroche, une phrase par groupe de documents, la problématique : trois mouvements suffisent. Les détails des documents viennent dans le développement.

\end{methode}

\hypertarget{uxe9tape-5-ruxe9diger-la-synthuxe8se-compluxe8te}{%
\subsection{Étape 5 --- Rédiger la synthèse complète}\label{uxe9tape-5-ruxe9diger-la-synthuxe8se-compluxe8te}}

\begin{indication}{1}

Ouvrez chaque partie par une phrase-thèse qui annonce l'idée, puis faites dialoguer les documents avec des connecteurs (\emph{whereas}, \emph{Kimball sees things differently}, \emph{Grovier replies}).

\end{indication}

\begin{indication}{2}

Citez peu et court, entre guillemets, et reformulez le reste : les citations ne doivent pas remplacer votre propre anglais.

\end{indication}

\hypertarget{solution-4}{%
\subsubsection{Solution}\label{solution-4}}

\textbf{Painting the Nation: Truth, Power and the Portrait}

Who decides how a nation sees itself? A 2001 photograph by David Dawson shows Lucian Freud painting Queen Elizabeth II, a portrait now held by London's National Portrait Gallery. Three articles, from The Guardian, the BBC and The Spectator, move the question to Washington, where the Smithsonian's museums have become a battlefield during Donald Trump's second term. They ask whether portraits and national symbols can still tell the truth once politics claims them.

The documents first agree that a portrait is a construction rather than a mirror. For Kim Sajet, director of the Smithsonian's National Portrait Gallery, ``all portraiture is a lie'': every image results from a negotiation between the sitter, the artist, the patron and an ever-changing audience. Dawson's photograph makes this visible. The Queen sits smiling and composed, wearing a diadem, while the painter, seen from behind, studies her before a small canvas on which her face looks heavier and sterner. National symbols are no different. Kelly Grovier recalls that the Statue of Liberty drew on Roman, Greek and Egyptian sources and has divided Americans since 1886, when suffragettes found it ironic and conservatives feared it would attract migrants. Amy Sherald's painting of Liberty as a black transgender woman reopens that old fault line.

Hence, the institutions that display such images have become political targets. Sajet's gallery strives to remain ``above the fray'': its labels try ``not to editorialise'' and rest on historical fact, as the carefully worded caption of Trump's photograph shows. Roger Kimball sees things differently. Applauding Trump's March executive order, ``Restoring Truth and Sanity to American History'', he accuses the Smithsonian of spreading a ``divisive, race-centered ideology'' and welcomes an official review of its exhibition texts and curatorial procedures, meant to ``celebrate American exceptionalism'' before the country's 250th anniversary. Grovier shows what such pressure means in practice, in the very museum Sajet presented as neutral: fearing for their funding, the organisers of Sherald's exhibition had second thoughts about her painting, and the artist cancelled the whole show, citing ``a culture of censorship''. Both camps thus claim to defend the truth, but they disagree on who should define it.

Finally, the documents question whether meaning can be controlled at all. Kimball urges conservatives to ``commandeer'' the institutions of elite culture, as the left has done since the 1960s. Grovier replies that symbols escape such control: excluding Sherald's work has only increased its impact, for nothing is ``more visible than something hidden''. Sajet trusts time and distance. Her gallery presents every president from George Washington to Barack Obama and Donald Trump, so that visitors can keep the present in perspective, and Trump's official portrait, though already completed, will stay in storage until 2029. As a Nigerian-born Dutchwoman, she also values an outsider's ability to stand back from America.

In the end, these documents suggest that no portrait tells the whole truth, yet, as Picasso put it, art may be ``a lie that makes us realise truth''. The real danger lies less in this partiality than in any attempt to impose a single, official image of a nation.

\textbf{Number of words (title included): 518}

\begin{methode}

Attribuez chaque position à son auteur (\emph{he accuses}, \emph{Grovier replies}, \emph{Sajet trusts}). La synthèse reste neutre, même quand un document est très engagé comme celui de Kimball.

\end{methode}

\begin{piege}

N'ajoutez aucun fait extérieur au dossier, même exact : ni ce que l'on sait par ailleurs de la réception du portrait de Freud, ni ce qui est arrivé ensuite aux personnes citées. Tout ce que dit la synthèse doit se retrouver dans les documents.

\end{piege}

\hypertarget{uxe9tape-6-compter-les-mots-et-relire}{%
\subsection{Étape 6 --- Compter les mots et relire}\label{uxe9tape-6-compter-les-mots-et-relire}}

\begin{indication}{1}

Comptez paragraphe par paragraphe au brouillon, titre compris, et reportez le total exact.

\end{indication}

\begin{indication}{2}

Relisez la langue : accords, temps (\emph{has divided \ldots{} since 1886}), articles, orthographe britannique ou américaine cohérente.

\end{indication}

\hypertarget{solution-5}{%
\subsubsection{Solution}\label{solution-5}}

\begin{longtable}[]{@{}ll@{}}
\toprule\noalign{}
Bloc & Mots \\
\midrule\noalign{}
\endhead
\bottomrule\noalign{}
\endlastfoot
Titre & 8 \\
Introduction & 73 \\
Partie I & 133 \\
Partie II & 147 \\
Partie III & 108 \\
Conclusion & 49 \\
\textbf{Total} & \textbf{518} \\
\end{longtable}

Le total (518 mots) est dans la fourchette 450-550, quelle que soit la convention adoptée pour les traits d'union et les apostrophes (voir la revue).

Points de langue vérifiés :

\begin{itemize}
\tightlist
\item
  \emph{has divided Americans since 1886} : \emph{present perfect} avec \emph{since} pour une situation qui dure encore ;
\item
  \emph{the left has done since the 1960s} : même règle ;
\item
  \emph{organisers}, \emph{editorialise}, \emph{realise} : orthographe britannique, reprise des documents et gardée d'un bout à l'autre ; \emph{race-centered} est une citation de Kimball et garde sa graphie américaine ;
\item
  \emph{had second thoughts about} : expression idiomatique reprise de Grovier et intégrée à la phrase.
\end{itemize}

\begin{piege}

Un nombre de mots inexact ou absent, ou un total hors fourchette, contrevient à une consigne explicite du sujet. Indiquer « about 500 words » ne répond pas à la demande d'indiquer le nombre « avec précision ».

\end{piege}

\hypertarget{bilan}{%
\section{Bilan}\label{bilan}}

\begin{itemize}
\tightlist
\item
  \textbf{Méthode de la synthèse} : analyse document par document, grille de confrontation, problématique tirée des désaccords, plan en trois parties qui croise les sources, rédaction neutre et attribuée.
\item
  \textbf{Consignes de forme} : titre obligatoire, décompte exact à la fin, fourchette de 450 à 550 mots, titre inclus.
\item
  \textbf{Document iconographique} : le décrire précisément (qui, quoi, comment) puis le relier à une idée du dossier.
\item
  \textbf{Langue utile} : \emph{to commission a portrait}, \emph{sitter}, \emph{patron}, \emph{curator}, \emph{label / caption}, \emph{to editorialise}, \emph{executive order}, \emph{to withdraw a work}, \emph{to cancel a show}, \emph{censorship}, \emph{culture wars}.
\item
  \textbf{Transposable} : les dossiers sur la liberté d'expression, la mémoire nationale, les musées et le pouvoir politique se traitent avec la même grille (neutralité revendiquée, intervention de l'État, résistance des symboles).
\end{itemize}

\end{document}
